\section{Introduction}

The expanding availability and the decreasing costs of use have made weather data, particularly remotely sensed data, a common input into socioeconomic analysis \citep{Dell2014-rf, DonaldsonStoreygard16}. Economists use these ready-made Earth observation products to predict a variety of economics outcomes, such as agricultural output, industrial output, labor productivity, energy demand, health, conflict, and economic growth. But, there are many ways to quantify weather and very little theoretical or systematic guidance on which weather metric is best suited to a given research question and geographic context. The results is that seemingly similar measures of ``rainfall,'' ``temperature,'' or ``shock'' can be based on very different summaries of the underlying weather data \citep{MichlerEtAl21WB}. The variety of metrics used and the size of weather data frames can create data handling and cleaning problem, meaning that the probability of making a coding error in calculating any given metric is non-trivial.

In this article, we introduce the community-contributed command \texttt{wxsum} which processes daily temperature and precipitation data and outputs useful summary statistics as well as measures of weather ``shocks.'' The goal of the command is to allow for easy and error free calculation of the most common weather metrics used in the socioeconomic literature. The \texttt{wxsum} command takes in daily temperature or precipitation data, typically in the form of gridded Earth observation data, and calculates a set of seasonal precipitation and temperature measures. These include a number of summary statistics for both temperature and precipitation, like the mean, median, standard deviation, and skew. It also calculates phenomenon-specific statistics, such as max seasonal temperature, growing degree days (GDD), and total seasonal rainfall. The command also calculates a number of measures of shocks. These include deviations from long-run averages, z-scores, dry spells, and killing degree days (KDD).

The \texttt{wxsum} command allows the user to set a variety of options so as to customize the weather metrics to their specific geography and research context. One can set the start and end day and month of a season, the bounds for GDD, the threshold of KDD, the number of temperature percentile bins, the threshold volume for what is considered a day with rain, and the length in years of what the user considers ``long run.''

The command itself is agnostic in terms of which of the many weather metrics that it calculates a researcher chooses to use. However, in creating the \texttt{wxsum} command, we are not endorsing a theory-free or ad hoc approach to incorporating measures of weather into socioeconomic analysis. There is a growing body of literature that raises concerns regarding the use (and misuse) of weather data, particularly when weather is used as a source of identifying exogenous variation \citep{BenamiEtAl26}. This literature raises concerns about the way gridded weather data is matched with socioeconomic data \citep{MichlerEtAl22}, the existence of measurement error in Earth observation weather data \cite{JosephsonEtAl26}, and the validity of the assumption that the weather is truly exogenous \citep{Sarsons15, Mellon2023-un}. While incorporating weather data into one's analysis is now low cost, a naive or unexamined use of that data will continue to have a high cost in terms of the accuracy of one's analysis and results.


\section{Data considerations}

The \texttt{wxsum} command is designed for a common structure of remotely sensed or gridded weather products after they have been matched to locations in an analysis dataset. Each observation is a location, plot, household, firm, or other spatial unit, and each weather variable is a daily measurement for that unit. Thus the command begins after the spatial matching step: the user should first decide which weather product, temporal aggregation, and spatial linking rule are appropriate for the research design. These choices remain substantive, not merely computational, because estimates can be sensitive to the source of weather data and to the way gridded observations are linked to socioeconomic observations \citep{AuffhammerEtAl13, DonaldsonStoreygard16, MichlerEtAl22, JosephsonEtAl26}.

\subsection{Shape and naming convention}

Input data must be in wide form. Daily weather values are stored in separate numeric variables whose names share a common prefix and end with an eight-digit date in \texttt{yyyymmdd} format. For example, if the precipitation variables are named with the prefix \texttt{rf\_}, then precipitation on May 15, 1979, should appear in a variable named \texttt{rf\_19790515}. If temperature variables are named with the prefix \texttt{tmp\_}, then temperature on the same date should appear in \texttt{tmp\_19790515}. Variables not beginning with the specified prefix are ignored unless they are retained by using the \texttt{keep()} option.

The wide representation is convenient for the gridded weather extracts commonly used in applied work because the unit of analysis stays in one row and each daily exposure is represented by a dated variable. It can also produce large datasets when daily histories cover many years or many spatial units. Thus the practical constraint is Stata's ability to hold the matched weather extract in memory rather than an additional conceptual limit imposed by \texttt{wxsum}.

This naming convention makes the calendar explicit. The command parses the date suffix of each weather variable, identifies the days that belong to each user-defined season, and computes statistics separately for each location and season. Because all computations are based on the dates in variable names, users should check that the date suffixes are complete, unique, and measured on the intended calendar. The command can accommodate seasons that cross calendar years. For example, a season that begins in November and ends in February is assigned to the year in which the season begins.

\subsection{Seasons, units, and missing values}

A season is defined by an initial month and final month and, optionally, by an initial day and final day. This flexibility is useful because the relevant period for rainfall or heat exposure differs across crops, agroecological zones, and outcomes. A growing season may be May through October in one application, whereas a heat-exposure season may begin late in the prior calendar year in another. The command does not impose a biological or econometric interpretation on the chosen season; the user should define the exposure window before running the command and should report that choice in subsequent empirical work.

The command is similarly agnostic about units. Rainfall can be measured in millimeters, centimeters, or inches, and temperature can be measured in degrees Celsius or Fahrenheit, provided that all daily variables are measured consistently. Unit choices matter for thresholds. The lower and upper bounds for growing degree days and the base for killing degree days must be expressed in the same units as the daily temperature data. The rainy-day threshold must be expressed in the same units as the daily precipitation data.

Many weather products contain missing values because of sensor problems, spatial coverage, cloud contamination, or limitations in station interpolation. Before using \texttt{wxsum}, users should examine missingness in the daily variables and decide whether the weather product and temporal coverage are adequate for the application. For means, medians, standard deviations, maxima, rainy-day counts, rain-day shares, temperature-bin shares, and dry spells, missing daily values are not treated as observed weather. For seasonal and monthly precipitation totals, \texttt{wxsum} uses Stata's \texttt{rowtotal()} function, so missing daily values contribute zero to those sums. If an empirical design requires a balanced panel of weather exposures or a minimum number of observed days within a season, that restriction should be imposed before or after running \texttt{wxsum} and should be documented in the replication files.


\section{Formulas and Calculations}

This section writes the variables generated by \texttt{wxsum} in notation. Let $i$ index the observation in the dataset and let $t$ denote the year in which a season begins. Let $S_t$ be the set of daily dates included in the season beginning in year $t$, and let $N_{it}$ be the number of nonmissing daily observations for observation $i$ in that season. For rainfall variables, let $R_{id}$ denote daily precipitation on date $d$; for temperature variables, let $T_{id}$ denote daily temperature. The command creates variables with the year suffix $\_t$, so the formulas below describe, for example, \texttt{mean\_}$t$, \texttt{total\_}$t$, and \texttt{gdd\_}$t$.

\subsection{Statistics common to rainfall and temperature}

For either rainfall or temperature data, write $X_{id}$ for the relevant daily variable, either $R_{id}$ or $T_{id}$. The daily mean is
\[
\texttt{mean\_}t_i = \frac{1}{N_{it}}\sum_{d\in S_t:X_{id}\ne .} X_{id} .
\]
The daily median, \texttt{median\_}$t$, is the median of the nonmissing values $\{X_{id}:d\in S_t, X_{id}\ne .\}$. The daily standard deviation is the sample standard deviation of those same nonmissing values,
\[
\texttt{sd\_}t_i = \left\{ \frac{1}{N_{it}-1}\sum_{d\in S_t:X_{id}\ne .}\left(X_{id}-\texttt{mean\_}t_i\right)^2 \right\}^{1/2} .
\]
The skewness variable is a standardized mean--median difference,
\[
\texttt{skew\_}t_i = \frac{\texttt{mean\_}t_i-\texttt{median\_}t_i}{\texttt{sd\_}t_i} .
\]
When the standard deviation is missing or zero, the skewness variable is missing.

\subsection{Rainfall calculations}

For precipitation data, the seasonal total is
\[
\texttt{total\_}t_i = \sum_{d\in S_t} R_{id},
\]
where the implementation follows \texttt{rowtotal()} and therefore treats missing daily rainfall components as zero in the sum. Monthly totals are computed first for each calendar month $m$ represented in the selected season,
\[
M_{imt}=\sum_{d\in S_t:month(d)=m} R_{id} .
\]
The variables \texttt{mean\_mo\_total\_}$t$, \texttt{median\_mo\_total\_}$t$, and \texttt{sd\_mo\_total\_}$t$ are respectively the mean, median, and sample standard deviation of the monthly totals $M_{imt}$ across the months represented in the season. The monthly skewness variable is
\[
\texttt{skew\_mo\_total\_}t_i = \frac{\texttt{mean\_mo\_total\_}t_i-\texttt{median\_mo\_total\_}t_i}{\texttt{sd\_mo\_total\_}t_i} .
\]
Let $c$ denote the value supplied in \texttt{rain\_threshold()}. Rainy days and days without rain are counted only among nonmissing daily rainfall values:
\[
\texttt{raindays\_}t_i = \sum_{d\in S_t:R_{id}\ne .} 1\{R_{id}\geq c\},
\]
\[
\texttt{norain\_}t_i = \sum_{d\in S_t:R_{id}\ne .} 1\{R_{id}<c\}.
\]
The share of observed days with rain is
\[
\texttt{pct\_raindays\_}t_i=\frac{\texttt{raindays\_}t_i}{N_{it}}.
\]
The dry-spell variable \texttt{dry\_}$t$ is the maximum length of a run of consecutive nonmissing days in the season for which $R_{id}<c$. Missing rainfall values break a dry spell rather than extending it.

\subsection{Temperature calculations}

For temperature data, the maximum is
\[
\texttt{max\_}t_i=\max_{d\in S_t:T_{id}\ne .} T_{id}.
\]
Let $a$ be the lower bound supplied in \texttt{gdd\_lo()} and $b$ be the upper bound supplied in \texttt{gdd\_hi()}. Daily growing degree days are truncated below at zero and above at $b-a$, and seasonal GDD are accumulated as
\[
\texttt{gdd\_}t_i = \sum_{d\in S_t:T_{id}\ne .} \min\left\{\max\left(T_{id}-a,0\right), b-a\right\} .
\]
If the user supplies a positive value $k$ in \texttt{kdd\_base()}, killing degree days are
\[
\texttt{kdd\_}t_i = \sum_{d\in S_t:T_{id}\ne .} \max\left(T_{id}-k,0\right) .
\]
If \texttt{kdd\_base()} is left at its default value of zero, \texttt{wxsum} does not generate the KDD variables.

The temperature-bin variables are season-specific shares. Let $B$ denote the number of bins requested in \texttt{bins()}, and let $q_{ibt}$ be the $100b/B$ percentile of the nonmissing daily temperatures in season $t$ for observation $i$, for $b=1,\ldots,B-1$. The first and last bins are
\[
\texttt{tempbin01\_}t_i=\frac{1}{N_{it}}\sum_{d\in S_t:T_{id}\ne .}1\{T_{id}<q_{i1t}\},
\]
\[
\texttt{tempbin}B\texttt{\_}t_i=\frac{1}{N_{it}}\sum_{d\in S_t:T_{id}\ne .}1\{T_{id}\geq q_{i,B-1,t}\}.
\]
For interior bins $b=2,\ldots,B-1$,
\[
\texttt{tempbin}b\texttt{\_}t_i=\frac{1}{N_{it}}\sum_{d\in S_t:T_{id}\ne .}1\{q_{i,b-1,t}\leq T_{id}<q_{ibt}\}.
\]
The generated bin numbers are zero-padded in the variable names, for example \texttt{tempbin01\_1993}. Across all detected seasons, \texttt{wxsum} also computes \texttt{binmean\_}$b$ as the mean of \texttt{tempbin}$b\_t$ over generated seasons and \texttt{binsd\_}$b$ as the corresponding sample standard deviation.

\subsection{Long-run reference periods}

Several variables generated by \texttt{wxsum} describe weather shocks as deviations from a location-specific long-run reference period. Let $w_{it}$ be one of the seasonal weather statistics for which a long-run benchmark is generated: \texttt{total\_}$t$, \texttt{raindays\_}$t$, \texttt{norain\_}$t$, or \texttt{pct\_raindays\_}$t$ for rainfall data; \texttt{gdd\_}$t$ and, when generated, \texttt{kdd\_}$t$ for temperature data. Let $L$ be the number of prior years requested in \texttt{lr\_years()}. The rolling long-run mean and standard deviation used for deviations and $z$ scores are based only on strictly preceding seasons:
\[
\bar{w}_{it}^{L}=L^{-1}\sum_{l=1}^{L} w_{i,t-l}
\]
\[
s_{it}^{L}=\left\{(L-1)^{-1}\sum_{l=1}^{L}\left(w_{i,t-l}-\bar{w}_{it}^{L}\right)^2\right\}^{1/2}.
\]
The deviation and corresponding $z$ score are
\[
\texttt{dev\_}w\texttt{\_}t_i=w_{it}-\bar{w}_{it}^{L},
\]
\[
\texttt{z\_}w\texttt{\_}t_i=\frac{w_{it}-\bar{w}_{it}^{L}}{s_{it}^{L}}.
\]
The use of strictly prior years avoids including the current realization in its own benchmark. If there are too few preceding season variables to satisfy the requested value of \texttt{lr\_years()}, the long-run deviations and $z$ scores are not generated for those initial seasons, although the nonreference statistics are still calculated.


\section{The \texttt{wxsum} command}

\subsection{Syntax}

The command is installed from the development repository by typing

\begin{stlog}
. local repo https://raw.githubusercontent.com/jdavidm/wxsum/master/
. net install wxsum, from("`repo'") replace
\end{stlog}

The syntax is

\begin{stsyntax}
wxsum {\it prefix\/} , ini\_month({\it month\/}) fin\_month({\it month\/})
    [ ini\_day({\it day\/}) fin\_day({\it day\/}) rain\_data temp\_data
      gdd\_lo({\it num\/}) gdd\_hi({\it num\/}) kdd\_base({\it num\/})
      bins({\it num\/}) lr\_years({\it num\/}) keep({\it varlist\/})
      save({\it filename\/}) rain\_threshold({\it num\/}) ]
\end{stsyntax}

Here \textit{prefix} is the common stub for the daily weather variables. It is not a varlist. If the data contain variables named \texttt{rf\_19790101}, \texttt{rf\_19790102}, and so on, the prefix supplied to the command is \texttt{rf\_}. Exactly one of \texttt{rain\_data} and \texttt{temp\_data} should be specified.

\subsection{Required variables}

The minimal dataset contains an identifier for each location and a set of daily weather variables. The identifier is not an argument of \texttt{wxsum}; it is retained with \texttt{keep()} if it should appear in the output dataset. The weather variables must be numeric and must use the naming convention described above. Table~\ref{tab:datareq} summarizes these requirements.

\begin{table}[htbp]
\caption{Data requirements for \texttt{wxsum}}
\label{tab:datareq}
\begin{center}
\begin{tabularx}{\textwidth}{lX}
\hline
Element & Requirement \\
\hline
Observation & One row per spatial or analytic unit. \\
Weather variables & One variable per day, with a common prefix and a \texttt{yyyymmdd} date suffix. \\
Date information & Encoded in the weather-variable names rather than stored in separate year, month, and day variables. \\
Rainfall data & Numeric daily precipitation variables; use \texttt{rain\_data}. \\
Temperature data & Numeric daily temperature variables; use \texttt{temp\_data}. \\
Other variables & Retain with \texttt{keep()} if identifiers, coordinates, or covariates should remain in the output. \\
\hline
\end{tabularx}
\end{center}
\end{table}

\subsection{Options}

\begin{description}
\item[\texttt{ini\_month(month)}] specifies the initial month of the season. This option is required.

\item[\texttt{fin\_month(month)}] specifies the final month of the season. This option is required. If the final month precedes the initial month in the calendar year, the command treats the season as crossing calendar years.

\item[\texttt{ini\_day(day)}] specifies the first day of the initial month to be included in the season. The default is \texttt{01}.

\item[\texttt{fin\_day(day)}] specifies the last day of the final month to be included in the season. The default is \texttt{01}.

\item[\texttt{rain\_data}] specifies that the daily variables contain precipitation. This option is mutually exclusive with \texttt{temp\_data}.

\item[\texttt{temp\_data}] specifies that the daily variables contain temperature. This option is mutually exclusive with \texttt{rain\_data}.

\item[\texttt{gdd\_lo(num)}] specifies the lower bound for growing degree days. This option is required when \texttt{temp\_data} is specified.

\item[\texttt{gdd\_hi(num)}] specifies the upper bound for growing degree days. This option is required when \texttt{temp\_data} is specified.

\item[\texttt{kdd\_base(num)}] specifies the threshold above which killing degree days are calculated. The threshold should be in the same unit as the daily temperature variables.

\item[\texttt{bins(num)}] specifies the number of equal-sized temperature percentile bins to create. The default is 4, the minimum is 4, and the maximum is 10.

\item[\texttt{lr\_years(num)}] specifies the number of strictly preceding years used to compute rolling deviations and $z$ scores. The default is 10, and the maximum is 50.

\item[\texttt{keep(varlist)}] retains variables from the input dataset in the output dataset. This option is typically used for identifiers, administrative codes, coordinates, or other merge keys.

\item[\texttt{save(filename)}] saves the resulting dataset to \textit{filename}. If \texttt{save()} is not specified, the generated variables remain in memory.

\item[\texttt{rain\_threshold(num)}] specifies the threshold used to define a rainy day. The default is 1.
\end{description}

\subsection{Output}

\texttt{wxsum} is a data-management command rather than an estimation command. Its primary output is a dataset containing the retained variables requested in \texttt{keep()} and the generated seasonal weather variables. Because the command generates one set of variables for each season in the data, the number of output variables increases with the time span of the daily weather data. Generated variables are indexed by the season-start year. For example, if a season begins in November 2001 and ends in February 2002, the resulting seasonal statistics are associated with 2001, not 2002.

When \texttt{rain\_data} is specified, the command produces the precipitation statistics in table~\ref{tab:rainout}. These variables combine simple summaries of the daily distribution, summaries of monthly totals, seasonal totals, rainy-day counts, and shock measures relative to the rolling long-run reference period.

\begin{table}[htbp]
\caption{Rainfall statistics generated by \texttt{wxsum}}
\label{tab:rainout}
\begin{center}
\begin{tabularx}{\textwidth}{lX}
\hline
Group & Generated statistics \\
\hline
Daily distribution & Mean, median, standard deviation, and skewness of daily precipitation within the season. \\
Monthly totals & Mean, median, standard deviation, and skewness of monthly precipitation totals within the season. \\
Seasonal total & Total precipitation in the season, deviation from the long-run average of total precipitation, and $z$ score of total precipitation. \\
Rainy days & Number of rainy days, number of days without rain, deviations from the corresponding long-run averages, and corresponding $z$ scores. \\
Rainy-day share & Share of observed days with rain, deviation from the long-run average share, and corresponding $z$ score. \\
Dry spells & Longest intra-seasonal dry spell. \\
\hline
\end{tabularx}
\end{center}
\end{table}

When \texttt{temp\_data} is specified, the command produces the temperature statistics in table~\ref{tab:tempout}. The growing degree day calculation uses \texttt{gdd\_lo()} and \texttt{gdd\_hi()} as lower and upper bounds, whereas killing degree days are calculated relative to \texttt{kdd\_base()} when a positive value of \texttt{kdd\_base()} is specified. Temperature bins report the share of observed seasonal days falling into the requested number of equal-sized percentile bins.

\begin{table}[htbp]
\caption{Temperature statistics generated by \texttt{wxsum}}
\label{tab:tempout}
\begin{center}
\begin{tabularx}{\textwidth}{lX}
\hline
Group & Generated statistics \\
\hline
Daily distribution & Mean, median, standard deviation, skewness, and maximum daily temperature within the season. \\
Growing degree days & Growing degree days, deviation from the long-run average of growing degree days, and $z$ score of growing degree days. \\
Killing degree days & Killing degree days, deviation from the long-run average of killing degree days, and $z$ score of killing degree days when \texttt{kdd\_base()} is positive. \\
Temperature bins & Share of observed days falling in each requested percentile bin, plus cross-season means and standard deviations of those bin shares. \\
\hline
\end{tabularx}
\end{center}
\end{table}

The command calculates the nonreference variables whenever the selected daily variables and season are available. Rolling deviations and $z$ scores require sufficient prior history at the same location. Consequently, if the data begin in 2000 and the user requests \texttt{lr\_years(10)}, the first eligible long-run deviations are not available until the season beginning in 2010. This behavior is intentional: the current season is never included in the reference period against which it is compared.

\section{Examples}

The repository includes small rainfall and temperature datasets that can be used to check installation and illustrate the syntax. These files are also useful as templates for checking the required wide layout and variable-naming convention. The examples below begin by loading those datasets, as they should in a reproducible do-file.

\subsection{Rainfall summaries}

Suppose that the daily rainfall variables in \texttt{rain.dta} are named with the prefix \texttt{rf\_}. The following commands calculate seasonal rainfall statistics for a May 15--October 15 growing season and save the resulting dataset:

\begin{stlog}
. use rain.dta, clear
. wxsum rf_, ini_month(05) fin_month(10) ini_day(15) ///
>     fin_day(15) rain_data lr_years(10) rain_threshold(1) ///
>     save(rainfall_stats.dta)
\end{stlog}

The output dataset contains one observation for each input location. For each season observed in the daily rainfall variables, \texttt{wxsum} adds the rainfall statistics listed in table~\ref{tab:rainout}. If the user also needs location identifiers or other merge variables, these should be supplied in \texttt{keep()}:

\begin{stlog}
. use rain.dta, clear
. wxsum rf_, ini_month(05) fin_month(10) ini_day(15) ///
>     fin_day(15) rain_data keep(id region) ///
>     save(rainfall_stats.dta)
\end{stlog}

The resulting file can then be merged to a socioeconomic dataset by the retained identifier variables and season-specific generated weather variables.

\subsection{Temperature summaries}

For temperature data, the user specifies \texttt{temp\_data} and supplies the growing-degree-day bounds. The next example uses a season that begins in November and ends in February. Because the final month precedes the initial month, the command treats the exposure window as spanning calendar years and assigns the season to the year in which it begins.

\begin{stlog}
. use temp.dta, clear
. wxsum tmp_, ini_month(11) fin_month(02) temp_data ///
>     gdd_lo(8) gdd_hi(32) kdd_base(32) bins(4) ///
>     keep(id region)
\end{stlog}

This command generates the temperature summaries listed in table~\ref{tab:tempout}. The values \texttt{gdd\_lo(8)}, \texttt{gdd\_hi(32)}, and \texttt{kdd\_base(32)} are illustrative; they should be chosen to match the units of the underlying temperature data and the biological or behavioral process being studied.

\subsection{Changing the long-run benchmark}

The default long-run reference period is 10 prior years. Researchers working with longer weather histories may prefer a longer benchmark, whereas applications with short panels may require a shorter one. The following command uses the prior 20 years as the rolling reference period:

\begin{stlog}
. use rain.dta, clear
. wxsum rf_, ini_month(05) fin_month(10) rain_data ///
>     lr_years(20) keep(id)
\end{stlog}

Changing \texttt{lr\_years()} affects only the rolling deviations and $z$ scores. It does not affect the seasonal means, medians, standard deviations, skewness measures, totals, counts, shares, dry-spell measures, or temperature-bin shares.

\section{Conclusions}

\texttt{wxsum} reduces the amount of custom code needed to transform daily weather data into seasonal measures commonly used in socioeconomic analysis. It standardizes the construction of summary statistics, threshold-based heat measures, rainy-day counts, dry spells, and rolling weather-shock measures while leaving the substantive choices--data source, spatial match, season definition, units, and thresholds--with the researcher. This division of labor is useful: the command lowers the probability of programming errors in repetitive weather processing tasks, but it does not substitute for careful research design.

Future development could extend \texttt{wxsum} in several directions. First, additional output formats could be useful for users who prefer long datasets with one row per location-season rather than wide datasets with one variable per statistic-year. Second, diagnostic options that report missing days, season length, and the number of prior years available for each generated shock measure would make quality control easier in large applications. Third, the command could be expanded to additional weather and climate variables, such as humidity, wind speed, vapor pressure deficit, or solar radiation, when these variables are available in daily gridded products. These extensions would preserve the main purpose of the command: to make weather-data processing in Stata more transparent, reproducible, and less error prone.

\endinput
